\documentclass[10pt,a4paper]{amsart}
\usepackage{amsmath,amssymb,amsthm,graphicx}
\usepackage{verbatim}
\newenvironment{reference}{\comment}{\endcomment}
%\newenvironment{reference}{}{}
\newenvironment{slogan}{\comment}{\endcomment}
\newenvironment{history}{\comment}{\endcomment}

% For commutative diagrams we use Xy-pic
\usepackage[all]{xy}

% We use 2cell for 2-commutative diagrams.
\xyoption{2cell}
% Theorem environments.
%
\theoremstyle{plain}
\newtheorem{theorem}[subsection]{Theorem}
\newtheorem{proposition}[subsection]{Proposition}
\newtheorem{lemma}[subsection]{Lemma}

\theoremstyle{definition}
\newtheorem{definition}[subsection]{Definition}
\newtheorem{example}[subsection]{Example}
\newtheorem{exercise}[subsection]{Exercise}
\newtheorem{situation}[subsection]{Situation}

\theoremstyle{remark}
\newtheorem{remark}[subsection]{Remark}
\newtheorem{remarks}[subsection]{Remarks}

\numberwithin{equation}{subsection}

% Macros
%
\def\lim{\mathop{\mathrm{lim}}\nolimits}
\def\colim{\mathop{\mathrm{colim}}\nolimits}
\def\Spec{\mathop{\mathrm{Spec}}}
\def\Hom{\mathop{\mathrm{Hom}}\nolimits}
\def\Ext{\mathop{\mathrm{Ext}}\nolimits}
\def\SheafHom{\mathop{\mathcal{H}\!\mathit{om}}\nolimits}
\def\SheafExt{\mathop{\mathcal{E}\!\mathit{xt}}\nolimits}
\def\Sch{\mathit{Sch}}
\def\Mor{\mathop{\mathrm{Mor}}\nolimits}
\def\Ob{\mathop{\mathrm{Ob}}\nolimits}
\def\Sh{\mathop{\mathit{Sh}}\nolimits}
\def\NL{\mathop{N\!L}\nolimits}
\def\CH{\mathop{\mathrm{CH}}\nolimits}
\def\proetale{{pro\text{-}\acute{e}tale}}
\def\etale{{\acute{e}tale}}
\def\QCoh{\mathit{QCoh}}
\def\Ker{\mathop{\mathrm{Ker}}}
\def\Im{\mathop{\mathrm{Im}}}
\def\Coker{\mathop{\mathrm{Coker}}}
\def\Coim{\mathop{\mathrm{Coim}}}

% Boxtimes
%
\DeclareMathSymbol{\boxtimes}{\mathbin}{AMSa}{"02}

%
% Macros for moduli stacks/spaces
%
\def\QCohstack{\mathcal{QC}\!\mathit{oh}}
\def\Cohstack{\mathcal{C}\!\mathit{oh}}
\def\Spacesstack{\mathcal{S}\!\mathit{paces}}
\def\Quotfunctor{\mathrm{Quot}}
\def\Hilbfunctor{\mathrm{Hilb}}
\def\Curvesstack{\mathcal{C}\!\mathit{urves}}
\def\Polarizedstack{\mathcal{P}\!\mathit{olarized}}
\def\Complexesstack{\mathcal{C}\!\mathit{omplexes}}
% \Pic is the operator that assigns to X its picard group, usage \Pic(X)
% \Picardstack_{X/B} denotes the Picard stack of X over B
% \Picardfunctor_{X/B} denotes the Picard functor of X over B
\def\Pic{\mathop{\mathrm{Pic}}\nolimits}
\def\Picardstack{\mathcal{P}\!\mathit{ic}}
\def\Picardfunctor{\mathrm{Pic}}
\def\Deformationcategory{\mathcal{D}\!\mathit{ef}}
\usepackage{cleveref}
\crefname{originallemma}{original lemma}{original lemmas}
\Crefname{originallemma}{Original lemma}{Original lemmas}
\makeatletter
\newlabel{lemma-dualizing-schemes}{{0A89}{source}{}{sourcepointer-1}{}}
\newlabel{lemma-dualizing-schemes@cref}{{[originallemma][0][]0A89}{[1][0][]source}{}{}{}}
\makeatother
\crefname{originallemma}{original lemma}{original lemmas}
\Crefname{originallemma}{Original lemma}{Original lemmas}
\makeatletter
\newlabel{perfect-lemma-quasi-coherence-internal-hom}{{0A6H}{source}{}{sourcepointer-2}{}}
\newlabel{perfect-lemma-quasi-coherence-internal-hom@cref}{{[originallemma][0][]0A6H}{[1][0][]source}{}{}{}}
\makeatother
\crefname{originallemma}{original lemma}{original lemmas}
\Crefname{originallemma}{Original lemma}{Original lemmas}
\makeatletter
\newlabel{more-algebra-lemma-base-change-RHom}{{0A6A}{source}{}{sourcepointer-3}{}}
\newlabel{more-algebra-lemma-base-change-RHom@cref}{{[originallemma][0][]0A6A}{[1][0][]source}{}{}{}}
\makeatother
\crefname{originallemma}{original lemma}{original lemmas}
\Crefname{originallemma}{Original lemma}{Original lemmas}
\makeatletter
\newlabel{lemma-dimension-function-scheme}{{0AWF}{source}{}{sourcepointer-4}{}}
\newlabel{lemma-dimension-function-scheme@cref}{{[originallemma][0][]0AWF}{[1][0][]source}{}{}{}}
\makeatother
\crefname{originallemma}{original lemma}{original lemmas}
\Crefname{originallemma}{Original lemma}{Original lemmas}
\makeatletter
\newlabel{dualizing-lemma-sitting-in-degrees}{{0A7U}{source}{}{sourcepointer-5}{}}
\newlabel{dualizing-lemma-sitting-in-degrees@cref}{{[originallemma][0][]0A7U}{[1][0][]source}{}{}{}}
\makeatother
\title{The degrees in which the Ext sheaves of a coherent module against a dualizing complex sit, in terms of depth}
\author{The Stacks Project Authors}
\date{stacks-project:duality.tex:local-snapshot-20261012; GNU FDL 1.2 or later}
\begin{document}
\maketitle
\noindent\textit{Editorial note.} All delivered mathematical text is the author's own, copied from the pinned Stacks Project chapter file duality.tex; no statement or proof step was rewritten, reordered or removed.\par\medskip
\noindent\textit{Editorial note.} Preamble provenance: the pinned chapter file opens with an input of the shared preamble and carries no preamble text of its own, so the project's single Stacks-wide preamble staged at sources/stacks-duality/preamble.tex is excerpted verbatim under the recipe fixed by LANE\_BRIEF 22.1, as three spans, lines 8-13 (verbatim package and the reference/slogan/history environments), lines 14-18 (Xy-pic setup) and lines 154-217 (the native theorem declarations and mathematical macros). Nothing was added, removed or redefined; the wrapper class and the amsmath/amssymb/amsthm/graphicx base packages are supplied by the assembler.\par\medskip
\noindent\textit{Editorial note.} The counted statement and proof contain no cross reference, no citation and no included file, so the delivered body is self contained.\par\medskip
\noindent\textit{Editorial note.} Canonical location of the counted proposition: Stacks tag 0ECM (sources/stacks-duality/pinned-original-source.tex line 588).\par\medskip
\noindent\textit{Allowed standard prerequisites (editorial).}\begin{itemize}
\item Coherence of the Ext sheaves against a dualizing complex. Same chapter, Lemma labelled lemma-dualizing-schemes, cited at proof lines 631-634. Gives coherence of E\textasciicircum{}i.
\item Stalk computation of the Ext sheaves and the dimension function. Derived Categories of Schemes, Lemma labelled perfect-lemma-quasi-coherence-internal-hom, cited at proof lines 634-638, More on Algebra, Lemma labelled more-algebra-lemma-base-change-RHom part (3), cited at proof lines 638-640, and Same chapter, Lemma labelled lemma-dimension-function-scheme, cited at proof lines 642-644. Reduce the local statements to the algebra.
\item The algebraic sitting-in-degrees statement. Dualizing Complexes, Lemma labelled dualizing-lemma-sitting-in-degrees, cited at proof lines 644-646. Supplies parts (1), (2) and (3).
\end{itemize}\medskip
\ifdefined\hypertarget\hypertarget{sourcepointer-1}{}\fi
\noindent\textit{Source reference.} Original lemma 0A89: sources/stacks-duality/pinned-original-source.tex, source lines 472--472. Same-chapter pointer into the author environment lemma at source line 472; invoked at proof line 609 of the counted proof. It is retained only as a navigation pointer; every load-bearing step is inside the delivered span.\par\medskip
\ifdefined\hypertarget\hypertarget{sourcepointer-2}{}\fi
\noindent\textit{Source reference.} Original lemma 0A6H: sources/stacks-chow/tags.txt, source lines 12845--12845. Cross-chapter standard prerequisite located in the project tag list staged in the local snapshot; invoked at proof line 612 of the counted proof. It is retained only as a navigation pointer; every load-bearing step is inside the delivered span.\par\medskip
\ifdefined\hypertarget\hypertarget{sourcepointer-3}{}\fi
\noindent\textit{Source reference.} Original lemma 0A6A: sources/stacks-chow/tags.txt, source lines 12838--12838. Cross-chapter standard prerequisite located in the project tag list staged in the local snapshot; invoked at proof line 615 of the counted proof. It is retained only as a navigation pointer; every load-bearing step is inside the delivered span.\par\medskip
\ifdefined\hypertarget\hypertarget{sourcepointer-4}{}\fi
\noindent\textit{Source reference.} Original lemma 0AWF: sources/stacks-duality/pinned-original-source.tex, source lines 567--567. Same-chapter pointer into the author environment lemma at source line 567; invoked at proof line 625 of the counted proof. It is retained only as a navigation pointer; every load-bearing step is inside the delivered span.\par\medskip
\ifdefined\hypertarget\hypertarget{sourcepointer-5}{}\fi
\noindent\textit{Source reference.} Original lemma 0A7U: sources/stacks-chow/tags.txt, source lines 12892--12892. Cross-chapter standard prerequisite located in the project tag list staged in the local snapshot; invoked at proof line 627 of the counted proof. It is retained only as a navigation pointer; every load-bearing step is inside the delivered span.\par\medskip
\begin{lemma}
\label{lemma-sitting-in-degrees}
Let $X$ be a locally Noetherian scheme. Let $\omega_X^\bullet$
be a dualizing complex on $X$ with associated dimension function $\delta$.
Let $\mathcal{F}$ be a coherent $\mathcal{O}_X$-module. Set
$\mathcal{E}^i = \SheafExt^{-i}_{\mathcal{O}_X}(\mathcal{F}, \omega_X^\bullet)$.
Then $\mathcal{E}^i$ is a coherent $\mathcal{O}_X$-module and
for $x \in X$ we have
\begin{enumerate}
\item $\mathcal{E}^i_x$ is nonzero only for
$\delta(x) \leq i \leq \delta(x) + \dim(\text{Supp}(\mathcal{F}_x))$,
\item $\dim(\text{Supp}(\mathcal{E}^{i + \delta(x)}_x)) \leq i$,
\item $\text{depth}(\mathcal{F}_x)$ is the smallest integer
$i \geq 0$ such that $\mathcal{E}_x^{i + \delta(x)} \not = 0$, and
\item we have
$x \in \text{Supp}(\bigoplus_{j \leq i} \mathcal{E}^j)
\Leftrightarrow
\text{depth}_{\mathcal{O}_{X, x}}(\mathcal{F}_x) + \delta(x) \leq i$.
\end{enumerate}
\end{lemma}

\begin{proof}
Lemma \ref{lemma-dualizing-schemes} tells us that $\mathcal{E}^i$
is coherent. Choosing an affine neighbourhood of $x$ and using
Derived Categories of Schemes, Lemma
\ref{perfect-lemma-quasi-coherence-internal-hom}
and
More on Algebra, Lemma
\ref{more-algebra-lemma-base-change-RHom} part (3)
we have
$$
\mathcal{E}^i_x =
\SheafExt^{-i}_{\mathcal{O}_X}(\mathcal{F}, \omega_X^\bullet)_x =
\Ext^{-i}_{\mathcal{O}_{X, x}}(\mathcal{F}_x,
\omega_{X, x}^\bullet) =
\Ext^{\delta(x) - i}_{\mathcal{O}_{X, x}}(\mathcal{F}_x,
\omega_{X, x}^\bullet[-\delta(x)])
$$
By construction of $\delta$ in Lemma \ref{lemma-dimension-function-scheme}
this reduces parts (1), (2), and (3) to
Dualizing Complexes, Lemma \ref{dualizing-lemma-sitting-in-degrees}.
Part (4) is a formal consequence of (3) and (1).
\end{proof}

\end{document}
